\par\SourceLeaf{055}%
\(3^\mathrm{o}\) partes vero reliquae, quae implicabunt vel factorem talem \(e^3e'\), vel talem \(eee'e''\), omnino omittuntur.

Valores medios biquadratorum \(e^4,\allowbreak e'^4,\allowbreak e''^4\) etc. ipsis biquadratis \(m^4,\allowbreak m'^4,\allowbreak m''^4\) etc. proportionales supponemus (vid. art.\ 16), ita ut illi sint ad haec ut \(\nu^4\) ad \(\mu^4\), adeoque \(\nu^4\) denotet valorem medium biquadratorum observationum talium quarum pondus\(=1\). Hinc praecepta praecedentia ita quoque exprimi poterunt: Loco singulorum biquadratorum \(\varepsilon^4,\varepsilon'^4,\varepsilon''^4\) etc. scribendum erit \(\nu^4\), loco singulorum productorum e binis quadratis ut \(\varepsilon\varepsilon\varepsilon'\varepsilon',\allowbreak\varepsilon\varepsilon\varepsilon''\varepsilon'',\allowbreak\varepsilon'\varepsilon'\varepsilon''\varepsilon''\) etc., scribendum erit \(\mu^4\), omnesque reliqui termini, qui implicabunt factores tales ut \(\varepsilon^3\varepsilon'\), vel \(\varepsilon\varepsilon\varepsilon'\varepsilon''\), vel \(\varepsilon\varepsilon'\varepsilon''\varepsilon'''\) erunt supprimendi.

His probe intellectis facile patebit

I. Valorem medium quadrati \(\Omega^0\Omega^0\) esse \(\pi\nu^4+(\pi\pi-\pi)\mu^4\)

II. Valor medius producti \(\varepsilon\varepsilon x^0\xi^0\) fit \(=a\alpha\nu^4+(a'\alpha'+a''\alpha''+\text{etc.})\mu^4\), sive quoniam \(a\alpha+a'\alpha'+a''\alpha''+\text{etc.}=1\),
\[
=a\alpha(\nu^4-\mu^4)+\mu^4
\]
Et quum perinde valor medius producti \(\varepsilon'\varepsilon'x^0\xi^0\) fiat \(=a'\alpha'(\nu^4-\mu^4)+\mu^4\), valor medius producti \(\varepsilon''\varepsilon''x^0\xi^0\) autem \(=a''\alpha''(\nu^4-\mu^4)+\mu^4\) et sic porro, patet, valorem medium producti \((\varepsilon\varepsilon+\varepsilon'\varepsilon'+\varepsilon''\varepsilon''+\text{etc.})x^0\xi^0\) sive \(\Omega^0x^0\xi^0\) esse
\[
=\nu^4-\mu^4+\pi\mu^4
\]
\begingroup\setlength{\rightskip}{0pt plus 2em}\noindent Eundem valorem medium habebunt producta \(\Omega^0y^0\eta^0,\allowbreak\Omega^0z^0\zeta^0\) etc. Quapropter valor medius producti \(\Omega^0(x^0\xi^0+y^0\eta^0+z^0\zeta^0+\text{etc.})\) fit\par\endgroup
\[
=\rho\nu^4+\rho(\pi-1)\mu^4
\]

III. Ne evolutiones reliquae nimis prolixae evadant, idonea denotatio introducenda erit. Utemur itaque characteristica \(\Sigma\) sensu aliquantum latiore quam supra passim factum est, ita ut denotet aggregatum termini, cui praefixa est, cum omnibus similibus sed non identicis inde per omnes observationum permutationes oriundis. Hoc pacto e.\ g. habemus \(x^0=\Sigma\alpha\varepsilon\), \(x^0x^0=\Sigma\alpha\alpha\varepsilon\varepsilon+2\Sigma\alpha\alpha'\varepsilon\varepsilon'\). Colligendo itaque valorem medium producti \(x^0x^0\xi^0\xi^0\) per partes, habemus primo valorem medium producti \(\alpha\alpha\varepsilon\varepsilon\xi^0\xi^0\)
\[
\begin{aligned}
&=aa\alpha\alpha\nu^4+\alpha\alpha(a'a'+a''a''+\text{etc.})\mu^4\\
&=aa\alpha\alpha(\nu^4-\mu^4)+\alpha\alpha\mu^4\Sigma aa
\end{aligned}
\]
\SourceLeaf{056}
Perinde valor medius producti \(\alpha'\alpha'\varepsilon'\varepsilon'\xi^0\xi^0\) fit \(=a'a'\alpha'\alpha'(\nu^4-\mu^4)+\alpha'\alpha'\mu^4\Sigma aa\) et sic porro, adeoque valor medius producti \(\xi^0\xi^0\Sigma\alpha\alpha\varepsilon\varepsilon\)
\[
=(\nu^4-\mu^4)\Sigma aa\alpha\alpha+\mu^4\Sigma aa.\Sigma\alpha\alpha
\]
Porro valor medius producti \(\alpha\alpha'\varepsilon\varepsilon'\xi^0\xi^0\) fit \(=2aa'\alpha\alpha'\mu^4\), valor medius producti \(\alpha\alpha''\varepsilon\varepsilon''\xi^0\xi^0\) perinde \(=2aa''\alpha\alpha''\mu^4\) etc., unde facile concluditur, valorem medium producti \(\xi^0\xi^0\Sigma\alpha\alpha'\varepsilon\varepsilon'\) fieri
\[
\begin{aligned}
&=2\mu^4\Sigma a\alpha a'\alpha'=\mu^4((\Sigma a\alpha)^2-\Sigma aa\alpha\alpha)\\
&=\mu^4(1-\Sigma aa\alpha\alpha)
\end{aligned}
\]
His collectis habemus valorem medium producti \(x^0x^0\xi^0\xi^0\)
\[
=(\nu^4-3\mu^4)\Sigma aa\alpha\alpha+2\mu^4+\mu^4\Sigma aa.\Sigma\alpha\alpha
\]

IV. Haud absimili modo invenitur valor medius producti \(x^0y^0\xi^0\eta^0\)
\[
\begin{aligned}
&=\nu^4\Sigma ab\alpha\beta+\mu^4\Sigma a\alpha b'\beta'\\
&\quad+\mu^4\Sigma ab\alpha'\beta'+\mu^4\Sigma a\beta b'\alpha'
\end{aligned}
\]
Sed habetur
\[
\begin{aligned}
\Sigma a\alpha b'\beta'&=\Sigma a\alpha.\Sigma b\beta-\Sigma a\alpha b\beta\\
\Sigma ab\alpha'\beta'&=\Sigma ab.\Sigma\alpha\beta-\Sigma ab\alpha\beta\\
\Sigma a\beta b'\alpha'&=\Sigma a\beta.\Sigma b\alpha-\Sigma a\beta b\alpha
\end{aligned}
\]
unde valor ille medius fit, propter \(\Sigma a\alpha=1,\allowbreak\Sigma b\beta=1,\allowbreak\Sigma a\beta=0,\allowbreak\Sigma b\alpha=0\),
\[
=(\nu^4-3\mu^4)\Sigma ab\alpha\beta+\mu^4(1+\Sigma ab.\Sigma\alpha\beta)
\]

V. Quum prorsus eodem modo valor medius producti \(x^0z^0\xi^0\zeta^0\) fiat
\[
=(\nu^4-3\mu^4)\Sigma ac\alpha\gamma+\mu^4(1+\Sigma ac.\Sigma\alpha\gamma)
\]
et sic porro, additio valorem medium producti \(x^0\xi^0(x^0\xi^0+y^0\eta^0+z^0\zeta^0+\text{etc.})\) suppeditat
\[
\begin{aligned}
&=(\nu^4-3\mu^4)\Sigma(a\alpha(a\alpha+b\beta+c\gamma+\text{etc.}))+(\rho+1)\mu^4\\
&\quad+\mu^4(\Sigma aa.\Sigma\alpha\alpha+\Sigma ab.\Sigma\alpha\beta+\Sigma ac.\Sigma\alpha\gamma+\text{etc.})\\
&=(\nu^4-3\mu^4)\Sigma(a\alpha(a\alpha+b\beta+c\gamma+\text{etc.}))+(\rho+2)\mu^4
\end{aligned}
\]

VI. Prorsus eodem modo valor medius producti \(y^0\eta^0(x^0\xi^0+y^0\eta^0+z^0\zeta^0+\text{etc.})\) eruitur
\SourceLeaf{057}
\[
=(\nu^4-3\mu^4)\Sigma(b\beta(a\alpha+b\beta+c\gamma+\text{etc.}))+(\rho+2)\mu^4
\]
dein valor medius producti \(z^0\zeta^0(x^0\xi^0+y^0\eta^0+z^0\zeta^0+\text{etc.})\)
\[
=(\nu^4-3\mu^4)\Sigma(c\gamma(a\alpha+b\beta+c\gamma+\text{etc.}))+(\rho+2)\mu^4
\]
et sic porro. Hinc per additionem prodit valor medius quadrati \((x^0\xi^0+y^0\eta^0+z^0\zeta^0+\text{etc.})^2\)
\[
=(\nu^4-3\mu^4)\Sigma((a\alpha+b\beta+c\gamma+\text{etc.})^2)+(\rho\rho+2\rho)\mu^4
\]

VII. Omnibus tandem rite collectis eruitur
\[
\begin{aligned}
N&=(\pi-2\rho)\nu^4+(\pi\pi-\pi-2\pi\rho+4\rho+\rho\rho)\mu^4\\
&\quad+(\nu^4-3\mu^4)\Sigma((a\alpha+b\beta+c\gamma+\text{etc.})^2)\\
&=(\pi-\rho)(\nu^4-\mu^4)+(\pi-\rho)^2\mu^4\\
&\quad-(\nu^4-3\mu^4)[\rho-\Sigma((a\alpha+b\beta+c\gamma+\text{etc.})^2)]
\end{aligned}
\]
Error itaque medius in determinatione ipsius \(\mu\mu\) per formulam
\[
\mu\mu=\frac{M}{\pi-\rho}
\]
metuendus erit
\[
=\BareRoot{\left\{\frac{\nu^4-\mu^4}{\pi-\rho}-\frac{\nu^4-3\mu^4}{(\pi-\rho)^2}.[\rho-\Sigma((a\alpha+b\beta+c\gamma+\text{etc.})^2)]\right\}}
\]

\Article{40}
Quantitas \(\Sigma((a\alpha+b\beta+c\gamma+\text{etc.})^2)\), quae in expressionem modo inventam ingreditur, generaliter quidem ad formam simpliciorem reduci nequit: nihilominus duo limites assignari possunt, inter quos ipsius valor necessario iacere debet. \textit{Primo} scilicet e relationibus supra evolutis facile demonstratur, esse
\[
\begin{aligned}
&(a\alpha+b\beta+c\gamma+\text{etc.})^2+(a\alpha'+b\beta'+c\gamma'+\text{etc.})^2\\
&\quad+(a\alpha''+b\beta''+c\gamma''+\text{etc.})^2+\text{etc.}\\
&=a\alpha+b\beta+c\gamma+\text{etc.}
\end{aligned}
\]
unde concludimus, \(a\alpha+b\beta+c\gamma+\text{etc.}\) esse quantitatem positivam unitate minorem (saltem non maiorem). Idem valet de quantitate \(a'\alpha'+b'\beta'+c'\gamma'+\text{etc.}\), quippe cui aggregatum
\[
\begin{aligned}
&(a'\alpha+b'\beta+c'\gamma+\text{etc.})^2+(a'\alpha'+b'\beta'+c'\gamma'+\text{etc.})^2\\
&\quad+(a'\alpha''+b'\beta''+c'\gamma''\text{ etc.})^2+\text{etc.}
\end{aligned}
\]
aequale invenitur; ac perinde \(a''\alpha''+b''\beta''+c''\gamma''+\text{etc.}\) unitate minor erit, et sic
